\section{Расчётная задача 4}
Пусть~$\ee = (e_1, e_2, e_3)$ --- базис в~$\kk^3$.
Пусть
\[
	x =
	\ee
	\Vect{x^1 \\ x^2 \\ x^3}
	\in \kk^3
	,
\]
а также
\begin{align*}
	\phi_1(x) &=
	\ee
	\Vect{
		x^2 - x^3 \\
		x^1 \\
		x^1 + x^3
	}
	,
	&
	\phi_2(x) &=
	\ee
	\Vect{
		x^2 \\
		2x^3
		x^1
	}
	.
\end{align*}
Найти формулу, определяющую действие оператора
\[
	\psi(x) :=
	\bigl(
		(\phi_2^2 + \phi_1)
		\circ
		\phi_2
	\bigr)
	(x)
	.
\]

\begin{proofm}
	% Заметим, что задав отображения~$\phi_\nu$ в координатах, мы выбрали базис в~$\kk^3$. Назовём его~$\ee = (e_1, e_2, e_3)$.
	Заметим, что отображения~$\phi_\nu$ --- линейные. Чтобы написать их матрицы~$A_\nu$, можно воспользоваться двумя приёмами. Изложим каждый из них на примере одной функции.

	\paragraph{Способ 1.} Вспомним, что по определению~$i$-ый столбец матрицы оператора~$\phi$ в паре базисов --- это координаты вектора~$\phi(e_i)$ в нужном базисе. В нашем случае базис образа и прообраза один и тот же~($\ee$):
	\[
		\phi_1(\ee) = \ee\,A_1
		.
	\]

	Поэтому подставим в~$\phi_1$ вектор~$e_1$, который имеет в базисе~$\ee$ координаты~$(1, 0, 0)^\top$:
	\[
		\phi_1
		\left( 
			\ee
			\Vect{
				1 \\ 0 \\ 0
			}
		\right)
		=
		\ee
		\Vect{
			0 - 0 \\
			1 \\
			1 + 0
		}
		=
		\ee
		\Vect{
			0 \\ 1 \\ 1
		}
		.
	\]
	Это и будет первый столбец матрицы~$A_1$. Теперь аналогично найдём~$\phi_1(e_2)$ и $\phi_1(e_3)$:
	\begin{align*}
		\phi_1(e_2) &=
		\phi_1
		\left( 
			\ee
			\Vect{
				0 \\ 1 \\ 0
			}
		\right)
		=
		\ee
		\Vect{
			1 \\ 0 \\ 0
		},
		\\
		\phi_1(e_3) &=
		\phi_1
		\left( 
			\ee
			\Vect{
				0 \\ 0 \\ 1
			}
		\right)
		=
		\ee
		\Vect{
			-1 \\ 0 \\ 1
		}.
	\end{align*}
	Поэтому
	\[
		A_1 =
		\begin{pmatrix}
			0 & 1 & -1
			\\
			1 & 0 & 0
			\\
			1 & 0 & 1
		\end{pmatrix}
		.
	\]

	\paragraph{Способ 2.} Запишем теперь отображение~$\phi_2$ сразу в координатах. Оно заключается в умножении матрицы~$A_2$ слева на столбец координат:
	\[
		\phi_2(x) =
		A_2
		\Vect{ x^1 \\ x^2 \\ x^3}
		.
	\]
	Поэтому (поменяем местами левую и правую части равенства)
	\[
		\begin{pmatrix}
			* & * & * \\
			* & * & * \\
			* & * & *
		\end{pmatrix}
		\Vect{
			x^1 \\ x^2 \\ x^3
		}
		=
		\Vect{
			x^2 \\ 2x^3 \\ x^1
		}
		=
		\Vect{
			\bigl( \phi_2(x) \bigr)^1
			\\
			\bigl( \phi_2(x) \bigr)^2
			\\
			\bigl( \phi_2(x) \bigr)^3
		}
		.
	\]
	Чтобы получить элемент~$\bigl( \phi_2(x) \bigr)^\nu$, нужно умножить $\nu$-ую строку на столбец $(x^1, x^2, x^3)^\top$.

	Обратим внимание на первую строку матрицы~$A_2$ и назовём элементы этой строки $(*_1, *_2, *_3)$:
	\[
		\begin{pmatrix}
			*_1 & *_2 & *_3
			\\
			{}
			\\
			{}
		\end{pmatrix}
		\Vect{
			x^1 \\ x^2 \\ x^3
		}
		=
		\Vect{
			*_1 x^1 + *_2 x^2 + *_3 x^3
			\\
			{}
			\\
			{}
		}
		=
		\Vect{
			x^2
			\\
			{}
			\\
			{}
		}
		.
	\]
	Ясно, что $*_1 = 0$, $*_2 = 1$, $*_3 = 0$. Поэтому матрица~$A_2$ теперь выглядит так:
	\[
		A_2 =
		\begin{pmatrix}
			0 & 1 & 0 \\
			* & * & * \\
			* & * & *
		\end{pmatrix}
		.
	\]

	Рассматривая аналогичным образом равенства
	\begin{align*}
		\begin{pmatrix}
			\\
			*_1 & *_2 & *_3
			\\
			{}
		\end{pmatrix}
		\Vect{
			x^1 \\ x^2 \\ x^3
		}
		&=
		\Vect{
			\\
			*_1 x^1 + *_2 x^2 + *_3 x^3
			\\
			{}
		}
		=
		\Vect{
			\\
			2x^3
			\\
			{}
		}
		,
		\\
		\begin{pmatrix}
			\\
			\\
			*_1 & *_2 & *_3
		\end{pmatrix}
		\Vect{
			x^1 \\ x^2 \\ x^3
		}
		&=
		\Vect{
			\\
			\\
			*_1 x^1 + *_2 x^2 + *_3 x^3
		}
		=
		\Vect{
			\\
			\\
			x^1
		}
		,
	\end{align*}
	находим, что вторая строка матрицы~$A_2$ имеет вид
	\[
		\Vect{
			\\
			0 & 0 & 2
			\\
			{}
		},
	\]
	а третья ---
	\[
		\Vect{
			\\
			\\
			1 & 0 & 0
		}.
	\]
	Поэтому
	\[
		A_2 =
		\begin{pmatrix}
			0 & 1 & 0
			\\
			0 & 0 & 2
			\\
			1 & 0 & 0
		\end{pmatrix}
		.
	\]

	\paragraph{Нахождение матрицы составного отображения.}
	Известно, что сумме линейных отображений отвечает сумма их матриц, а их композиции (произведению) --- произведение матриц. Поэтому отображению
	\[
		\psi = (\phi_2^2 + \phi_1) \circ \phi_2
	\]
	отвечает матрица
	\[
		B := (A_2^2 + A_1) \cdot A_2
		.
	\]
	Можно раскрыть скобки:
	\[
		B = A_2^3 + A_1 A_2
		.
	\]
	Найдём каждое слагаемое отдельно. Для начала,
	\begin{align*}
		A_2^2 &=
		\left(
			\begin{array}{ccc}
				0 & 1 & 0
				\\
				\hline
				0 & 0 & 2
				\\
				\hline
				1 & 0 & 0
			\end{array}
		\right)
		\left(
		\begin{array}{c|c|c}
			0 & 1 & 0
			\\
			0 & 0 & 2
			\\
			1 & 0 & 0
		\end{array}
		\right)
		=
		\begin{pmatrix}
			0 & 0 & 2
			\\
			2 & 0 & 0
			\\
			0 & 1 & 0
		\end{pmatrix}
		;
		\\
		A_2^3 = A_2^2 \cdot A_2 &=
		\left(
			\begin{array}{ccc}
				0 & 0 & 2
				\\
				\hline
				2 & 0 & 0
				\\
				\hline
				0 & 1 & 0
			\end{array}
		\right)
		\left(
			\begin{array}{c|c|c}
				0 & 1 & 0
				\\
				0 & 0 & 2
				\\
				1 & 0 & 0
			\end{array}
		\right)
		=
		\begin{pmatrix}
			2 & 0 & 0
			\\
			0 & 2 & 0
			\\
			0 & 0 & 2
		\end{pmatrix}
		=
		2E
		,
	\end{align*}
	где $E$ --- единичная матрица. Но~$E$ соответствует тождественному оператору~$\id$; таким образом,~$\phi_2^3$ --- это оператор~$2\id$ растяжения в~$2$ раза.

	Теперь вычислим~$A_1 A_2$:
	\begin{align*}
		A_1 A_2 &=
		\left( 
			\begin{array}{rrr}
				0 & 1 & -1
				\\
				\hline
				1 & 0 & 0
				\\
				\hline
				1 & 0 & 1
			\end{array}
		\right)
		\left( 
			\begin{array}{c|c|c}
				0 & 1 & 0
				\\
				0 & 0 & 2
				\\
				1 & 0 & 0
			\end{array}
		\right)
		=
		\left( 
			\begin{array}{rrr}
				-1 & 0 & 2
				\\
				0 & 1 & 0
				\\
				1 & 1 & 0
			\end{array}
		\right)
		.
	\end{align*}
	Теперь осталось собрать всё вместе:
	\[
		B = A_2^3 + A_1 A_2 =
		2
		\begin{pmatrix}
			1 & 0 & 0
			\\
			0 & 1 & 0
			\\
			0 & 0 & 1
		\end{pmatrix}
		+
		\begin{pmatrix}
			-1 & 0 & 2
			\\
			0 & 1 & 0
			\\
			1 & 1 & 0
		\end{pmatrix}
		=
		\begin{pmatrix}
			1 & 0 & 2
			\\
			0 & 2 & 0
			\\
			1 & 1 & 2
		\end{pmatrix}
		.
	\]

	Теперь несложно понять, как действует~$\psi$ в координатах:
	\[
		\psi(x)
		=
		B
		\Vect{x^1 \\ x^2 \\ x^3}
		=
		\begin{pmatrix}
			1 & 0 & 2
			\\
			0 & 2 & 0
			\\
			1 & 1 & 2
		\end{pmatrix}
		\Vect{x^1 \\ x^2 \\ x^3}
		=
		\Vect{
			x^1 + 2x^3
			\\
			2x^2
			\\
			x^1 + x^2 + 2x^3
		}
		.
		\proofmark
	\]
\end{proofm}

